\documentclass[12pt]{article}
\usepackage[utf8]{inputenc}
\usepackage[T1]{fontenc}
\usepackage{hyperref}
\title{Iyasa Language Resources for K-12 Curriculum Development}
\author{Compiled by Claude Sonnet 5.5 (claude-sonnet-5-5) for Andr\'as Kornai}
\date{}
\begin{document}
\maketitle

\section{Introduction}
Iyasa (also Yasa or Yassa; ISO 639-3 yko, Glottocode yasa1242) is a Bantu language of the Atlantic-Congo family. It is spoken on the coast of Cameroon south of Kribi, including Campo, and across the Ntem River in Equatorial Guinea. Estimates are about 2,400 speakers (2000--2011) and about 3,000 (possibly an overestimate). It may be a dialect of Kombe, with which it is almost fully mutually intelligible. UNESCO classes it as definitely endangered, and younger speakers are shifting to French. An alphabet was adopted in Cameroon in 1979 and a new orthography was introduced in 2015. The search found no school materials.

\section{Primary Educational Resources}
The searches found no primers, textbooks or school programmes in Iyasa.

\section{Linguistic and Reference Materials}
\begin{itemize}
\item Wikipedia, ``Yasa language''. \url{https://en.wikipedia.org/wiki/Yasa_language}
\item Omniglot, Yasa: a page on the writing system. \url{https://omniglot.com/writing/yasa.htm}
\end{itemize}
I saw no grammar or dictionary in the results.

\section{Religious and Community Resources}
\begin{itemize}
\item Joshua Project, Yasa in Cameroon: people-group profile. \url{https://www.joshuaproject.net/index.php/people_groups/15998/CM}
\end{itemize}

\section{Digital and Audio Resources}
\begin{itemize}
\item Endangered Languages Project entry for Yasa (yko). \url{https://endangeredlanguages.com/lang/yko?hl=fr}
\end{itemize}
I found no audio or video resources.

\section{Curriculum Development Context}
The 2015 orthography is a starting point for written materials, but I found no evidence of it being used in schools. Because of the shift to French, any curriculum would need community-led work, and Kombe materials might be a useful comparison if the two are closely related. Documentation is thin, so new primers, wordlists and recordings would come first.

\end{document}